\documentclass[11pt]{article}

\usepackage[utf8]{inputenc}
\usepackage[T1]{fontenc}
\usepackage[spanish]{babel}
\usepackage{amsmath, amssymb}
\usepackage{geometry}
\usepackage{hyperref}
\usepackage{enumitem}
\usepackage{longtable}
\usepackage{array}

\geometry{margin=2.5cm}
\setlist[itemize]{topsep=4pt,itemsep=3pt}
\setlist[enumerate]{topsep=4pt,itemsep=3pt}

\title{Explicaci\'on detallada de RankGA}
\author{Documento t\'ecnico alineado con la implementaci\'on actual del repositorio}
\date{\today}

\begin{document}

\maketitle

\section{Prop\'osito de este documento}

Este documento explica con detalle el algoritmo implementado en el repositorio
\texttt{RankGA}. No describe un algoritmo gen\'erico inventado aparte del
c\'odigo, sino la versi\'on concreta que hoy existe en las clases
\texttt{rankga.RankGA}, \texttt{rankga.Population},
\texttt{rankga.Individual}, \texttt{rankga.Problem} y \texttt{rankga.Gene}.

La meta es dejar claro:

\begin{itemize}
  \item qu\'e problema resuelve el framework;
  \item cu\'al es la arquitectura conceptual del sistema;
  \item c\'omo opera cada fase de la evoluci\'on;
  \item qu\'e ecuaciones usa realmente la implementaci\'on;
  \item qu\'e informaci\'on genera como salida;
  \item y cu\'ales son sus supuestos, limitaciones y detalles que pueden
  provocar mala interpretaci\'on si solo se mira el nombre de las clases.
\end{itemize}

Este texto es deliberadamente exhaustivo. La idea no es solo decir ``qu\'e
hace'', sino dejar expl\'icito \emph{por qu\'e} est\'a organizado as\'i, qu\'e
convenciones adopta y d\'onde conviene tener cuidado.

\section{Idea general del algoritmo}

RankGA es un algoritmo gen\'etico con cuatro rasgos estructurales centrales:

\begin{enumerate}
  \item \textbf{Selecci\'on por rango.}
  La probabilidad esperada de supervivencia o replicaci\'on de un individuo no
  depende directamente de su valor num\'erico de fitness, sino de su
  \emph{posici\'on ordenada} dentro de la poblaci\'on.

  \item \textbf{Mutaci\'on dependiente del rango.}
  La intensidad de mutaci\'on no es uniforme en toda la poblaci\'on. Los
  individuos mejor rankeados reciben menos perturbaci\'on y los peor rankeados
  reciben m\'as.

  \item \textbf{Emparejamiento por rango para recombinaci\'on.}
  La recombinaci\'on no empareja individuos al azar. Despu\'es de ordenar la
  poblaci\'on por fitness, el algoritmo empareja individuos contiguos en ese
  orden. Por tanto, la estructura del rango no solo controla la selecci\'on y
  la mutaci\'on, sino tambi\'en \emph{qui\'en se recombina con qui\'en}.

  \item \textbf{Separaci\'on entre motor y problema.}
  El motor evolutivo es gen\'erico y delega en el problema concreto la forma de
  representar genes, construir individuos, evaluar fitness y definir el
  significado de la intensidad de mutaci\'on.
\end{enumerate}

La orientaci\'on del framework es de \textbf{maximizaci\'on}. Un fitness mayor
siempre es mejor. Toda la l\'ogica de ordenamiento y selecci\'on est\'a escrita
con esa convenci\'on.

\section{Arquitectura del sistema}

\subsection{Componentes principales}

El n\'ucleo del framework se reparte en estas abstracciones:

\begin{longtable}{|p{3cm}|p{10cm}|}
\hline
\textbf{Componente} & \textbf{Responsabilidad principal} \\
\hline
\texttt{Problem} & Define el problema de optimizaci\'on: fitness, longitud de
genoma, construcci\'on de genes e individuos, intensidad local, intensidad
global y fitness objetivo. \\
\hline
\texttt{Gene} & Define la unidad at\'omica del genoma. Cada implementaci\'on
decide c\'omo almacenar el valor, c\'omo mutarlo y c\'omo medir distancia a
otro gen hom\'ologo. \\
\hline
\texttt{Individual} & Representa una soluci\'on completa: genoma, fitness,
rango, intensidad de mutaci\'on recibida y operaciones de mutaci\'on y
recombinaci\'on a nivel del individuo. \\
\hline
\texttt{Population} & Encapsula las operaciones poblacionales: evaluaci\'on,
ordenamiento por fitness, selecci\'on, recombinaci\'on y mutaci\'on. \\
\hline
\texttt{RankGA} & Es el \emph{driver} de la corrida: inicializa, itera,
registra salidas, decide el paro y genera res\'umenes y gr\'aficas. \\
\hline
\texttt{AdaptiveProblem} & Interfaz opcional para problemas cuyo estado interno
puede adaptarse durante la corrida. \\
\hline
\texttt{ProblemFactory} & Selecci\'on expl\'icita del problema desde la l\'inea
de comandos. \\
\hline
\end{longtable}

\subsection{Separaci\'on entre motor y codificaci\'on}

El framework no impone una sola codificaci\'on:

\begin{itemize}
  \item puede haber genes enteros categ\'oricos;
  \item genes reales de precisi\'on doble;
  \item y cualquier otra representaci\'on que respete la interfaz
  \texttt{Gene}.
\end{itemize}

Esto tiene una consecuencia importante: la palabra \emph{mutaci\'on} no
significa lo mismo en todas las implementaciones. El motor solo calcula una
\textbf{intensidad} y la pasa a cada gen. El \texttt{Gene} concreto decide qu\'e
hacer con esa intensidad.

\section{Contrato formal entre problema y motor}

\subsection{La interfaz \texttt{Problem}}

Todo problema debe proveer:

\begin{itemize}
  \item \texttt{fitness(Individual)}: eval\'ua la calidad de una soluci\'on;
  \item \texttt{getProblemName()}: nombre visible del problema;
  \item \texttt{getGenomeLength()}: longitud del genoma;
  \item \texttt{getNewGene(boolean, Random)}: construcci\'on de genes nuevos;
  \item \texttt{getNewGene(Gene)}: copia o proyecci\'on de un gen existente;
  \item \texttt{getGoalFt()}: fitness objetivo que, si se alcanza, detiene la
  corrida;
  \item \texttt{getDisplayModulus()}: agrupaci\'on visual del genoma en logs;
  \item \texttt{getNewIndividual(boolean, Random)}: construcci\'on de
  individuos nuevos;
  \item \texttt{getNewIndividual(Individual)}: copia profunda de un individuo;
  \item \texttt{getLocalSearchIntensity()} y
  \texttt{getGlobalSearchIntensity()}.
\end{itemize}

\subsection{Interpretaci\'on de las intensidades \(L\) y \(G\)}

El problema define dos par\'ametros:

\begin{itemize}
  \item \(L\): intensidad local;
  \item \(G\): intensidad global.
\end{itemize}

En el framework actual:

\begin{itemize}
  \item \(G\) es la intensidad m\'axima asignable, t\'ipicamente asociada a los
  peores rangos;
  \item \(L\) se usa para calibrar la pendiente del perfil de mutaci\'on.
\end{itemize}

El supuesto matem\'atico deseado es \(G > L\), porque entonces el exponente de
mutaci\'on resulta positivo y la mutaci\'on crece con el rango. En la
implementaci\'on actual eso no se valida en tiempo de ejecuci\'on; simplemente
se asume.

\section{Representaci\'on de individuos y genes}

\subsection{Estructura de un individuo}

Cada individuo contiene:

\begin{itemize}
  \item un arreglo ordenado de genes;
  \item un fitness cacheado;
  \item un rango entero;
  \item la \'ultima intensidad de mutaci\'on aplicada;
  \item el rango de la pareja usada en la \'ultima recombinaci\'on;
  \item un campo textual auxiliar para logging.
\end{itemize}

La copia de un individuo es \textbf{profunda} a nivel de genes: el constructor
de copia llama, para cada locus, a \texttt{problem.getNewGene(gene)}.

\subsection{Importancia del contrato de copia}

Este detalle es m\'as importante de lo que parece. Si la copia de genes no se
hace correctamente, la independencia estoc\'astica entre individuos puede
romperse y la interpretaci\'on experimental del algoritmo deja de ser fiable.

La lecci\'on metodol\'ogica es clara:

\begin{itemize}
  \item el operador de copia no es un detalle administrativo;
  \item forma parte esencial de la sem\'antica del algoritmo;
  \item una copia incorrecta puede falsear conclusiones experimentales.
\end{itemize}

\subsection{Dos ejemplos de genes del repositorio}

\paragraph{GeneInteger.}
Representa un gen discreto en
\(\{0,1,\dots,\texttt{NUM\_VALUES}-1\}\). Su mutaci\'on interpreta la intensidad
como probabilidad de reemplazo por otro valor distinto del actual. Es decir, en
este caso la intensidad s\'i termina siendo usada como probabilidad, pero esa
decisi\'on se toma a nivel del gen, no del motor.

\paragraph{GeneDoublePrecision.}
Representa un gen real. Su mutaci\'on aplica ruido gaussiano aditivo con
desviaci\'on est\'andar igual a la intensidad recibida.

\section{Ciclo evolutivo completo}

\subsection{Esquema general}

Para un problema dado, una semilla, un tama\~no de poblaci\'on \(N\) y un
n\'umero de repeticiones \(R\), el ciclo externo es:

\begin{enumerate}
  \item crear una poblaci\'on inicial aleatoria;
  \item evaluarla y ordenarla;
  \item registrar un snapshot inicial;
  \item repetir mientras no se alcance el fitness objetivo y no se haya agotado
  la paciencia:
  \begin{enumerate}
    \item seleccionar;
    \item recombinar;
    \item reevaluar;
    \item registrar mejora si hubo;
    \item mutar;
    \item reevaluar;
    \item registrar mejora si hubo;
    \item adaptar el problema si implementa \texttt{AdaptiveProblem}.
  \end{enumerate}
  \item registrar un snapshot final;
  \item a\~nadir una fila al \texttt{summary.csv};
  \item al final de todas las repeticiones, generar gr\'aficas autom\'aticamente.
\end{enumerate}

\subsection{Inicializaci\'on}

La poblaci\'on inicial se crea con:

\begin{itemize}
  \item tama\~no fijo \(N\);
  \item individuos aleatorizados;
  \item un \texttt{Random} derivado de la semilla base de la corrida y del
  \'indice de repetici\'on.
\end{itemize}

Despu\'es se eval\'ua a toda la poblaci\'on y se ordena de fitness mayor a
menor. El n\'umero inicial de evaluaciones de fitness es exactamente \(N\).

\section{Selecci\'on por rango}

\subsection{Idea}

La selecci\'on no usa directamente el valor crudo del fitness. En su lugar:

\begin{enumerate}
  \item la poblaci\'on ya est\'a ordenada de mejor a peor;
  \item cada individuo recibe un rango impl\'icito por su posici\'on;
  \item ese rango determina cu\'antos clones se espera producir.
\end{enumerate}

\subsection{Presi\'on selectiva fija}

La implementaci\'on actual fija la presi\'on selectiva en
\[
S = 3.
\]

No es un par\'ametro configurable; es una decisi\'on de dise\~no del algoritmo.

\subsection{Definici\'on del rango}

En la selecci\'on, el c\'odigo usa
\[
r_i = \frac{i}{N-1},
\]
donde:

\begin{itemize}
  \item \(i=0\) corresponde al mejor individuo;
  \item \(i=N-1\) corresponde al peor;
  \item por lo tanto \(r_i \in [0,1]\).
\end{itemize}

En la implementaci\'on que este documento describe, la definici\'on correcta es
\(r_i = i/(N-1)\).

\subsection{Conteo esperado de clones}

El n\'umero esperado de clones del individuo \(i\) es
\[
c_i = S(1-r_i)^{S-1}.
\]

Como \(S=3\), se vuelve
\[
c_i = 3(1-r_i)^2.
\]

\subsection{Materializaci\'on en el c\'odigo}

La implementaci\'on hace dos pasos:

\begin{enumerate}
  \item asignaci\'on determinista:
  \[
  n_i = \lfloor c_i \rfloor;
  \]
  \item redondeo estoc\'astico sobre la parte fraccional hasta completar
  exactamente \(N\) individuos.
\end{enumerate}

Es decir, si
\[
f_i = c_i - \lfloor c_i \rfloor,
\]
entonces se hacen pasadas c\'iclicas y se agrega un clon extra al individuo
\(i\) con probabilidad \(f_i\), mientras el total de clones siga siendo menor
que \(N\).

\subsection{Observaciones t\'ecnicas}

\begin{itemize}
  \item La selecci\'on conserva exactamente el tama\~no de poblaci\'on.
  \item El procedimiento usa redondeo estoc\'astico secuencial porque eso basta
  para materializar la ley de clonaci\'on adoptada por el dise\~no.
  \item El redondeo estoc\'astico introduce un peque\~no sesgo de orden por el
  recorrido secuencial, pero actualmente ese sesgo se acepta como parte del
  dise\~no.
\end{itemize}

\section{Emparejamiento por rango y recombinaci\'on}

\subsection{Por qu\'e el emparejamiento por rango es central}

En este framework, el rango no solo decide cu\'antos clones sobreviven y qu\'e
intensidad de mutaci\'on recibe cada individuo. Tambi\'en decide la
\textbf{vecindad reproductiva}.

Eso significa que el rango organiza simult\'aneamente tres mecanismos:

\begin{enumerate}
  \item la presi\'on selectiva;
  \item la intensidad de mutaci\'on;
  \item y la estructura de apareamiento.
\end{enumerate}

Por eso el emparejamiento por rango debe considerarse una caracter\'istica
central del algoritmo y no un detalle de implementaci\'on. Si el emparejamiento
fuera aleatorio, la sem\'antica del m\'etodo cambiar\'ia de manera importante,
porque se perder\'ia la relaci\'on sistem\'atica entre orden por fitness y
mezcla gen\'etica.

\subsection{Emparejamiento}

La recombinaci\'on se hace en parejas adyacentes:
\[
(0,1), (2,3), (4,5), \dots
\]

No hay emparejamiento aleatorio global. El criterio es puramente posicional
despu\'es del ordenamiento por fitness.

\subsection{Operador}

El operador de recombinaci\'on a nivel de individuo es cruzamiento uniforme con
\textbf{intercambio in situ}. Este detalle es esencial:

\begin{itemize}
  \item la selecci\'on ya produjo una nueva poblaci\'on compuesta por clones;
  \item la recombinaci\'on no construye objetos ``hijo'' adicionales;
  \item en su lugar, toma dos clones ya existentes y les intercambia genes
  locus por locus.
\end{itemize}

Por tanto, la secuencia correcta del algoritmo es:

\begin{enumerate}
  \item la selecci\'on multiplica genotipos mediante clonaci\'on;
  \item la poblaci\'on clonada reemplaza a la anterior;
  \item la recombinaci\'on act\'ua sobre esa poblaci\'on ya clonada;
  \item cada pareja de clones se modifica directamente por intercambio de
  loci.
\end{enumerate}

No hay una fase separada de ``generaci\'on de descendencia'' en el sentido
cl\'asico de crear dos hijos nuevos a partir de dos padres intactos. Los
``padres'' que entran a recombinaci\'on son exactamente los individuos clonados
que sobreviven a la selecci\'on, y despu\'es del intercambio esos mismos
objetos quedan convertidos en las soluciones recombinadas que siguen al resto
del ciclo.

En cada pareja, el mecanismo es:

\begin{itemize}
  \item para cada locus \(j\),
  \item con probabilidad \(0.5\),
  \item se intercambian los genes de ambos individuos en ese locus.
\end{itemize}

Formalmente, si \(x_j\) y \(y_j\) son los genes de dos padres, entonces:

\[
(x_j, y_j) \leftarrow
\begin{cases}
(y_j, x_j), & \text{con probabilidad } 0.5, \\
(x_j, y_j), & \text{con probabilidad } 0.5.
\end{cases}
\]

\subsection{Consecuencias}

\begin{itemize}
  \item Los individuos cercanos en rango tienden a recombinarse entre s\'i.
  \item La mezcla gen\'etica est\'a sesgada por la estructura del ordenamiento.
  \item La recombinaci\'on opera sobre clones seleccionados, no sobre una
  poblaci\'on parental separada de una poblaci\'on filial.
  \item El operador no aumenta el tama\~no de la poblaci\'on, porque no crea
  individuos nuevos; solo transforma los ya existentes.
  \item El operador es sim\'etrico.
  \item No privilegia ni prefijos ni sufijos del genoma.
  \item Conserva la distribuci\'on marginal de genes, aunque mezcla bloques de
  construcci\'on a nivel de loci.
\end{itemize}

\section{Mutaci\'on por rango}

\subsection{Idea central}

La mutaci\'on no es plana. La intensi\'on del dise\~no es:

\begin{itemize}
  \item preservar m\'as a los individuos bien rankeados;
  \item perturbar m\'as a los individuos peor rankeados;
  \item hacer coexistir explotaci\'on cerca del frente y exploraci\'on en la
  cola de la poblaci\'on.
\end{itemize}

\subsection{Exponente de mutaci\'on}

El framework define
\[
\beta = \frac{\ln(G/L)}{\ln(N-1)}.
\]

Si \(G>L\), entonces \(\beta>0\).

\subsection{Intensidad por individuo}

Para el individuo de rango \(i\), la intensidad es
\[
I_i = G \, r_i^\beta,
\]
con
\[
r_i = \frac{i}{N-1}.
\]

De aqu\'i se sigue que:

\begin{itemize}
  \item para el mejor individuo, \(r_0=0\) y por tanto \(I_0=0\);
  \item para el segundo mejor individuo, \(r_1=\frac{1}{N-1}\), y entonces
  \[
  I_1 = G \left(\frac{1}{N-1}\right)^\beta = L;
  \]
  \item para el peor individuo, \(r_{N-1}=1\) y por tanto
  \(I_{N-1}=G\);
  \item la intensidad crece mon\'otonamente con el rango si \(\beta>0\).
\end{itemize}

Ese detalle no es accidental. La forma en que se define \(\beta\) hace que el
perfil de intensidades quede anclado exactamente en tres puntos:

\[
I_0 = 0,\qquad I_1 = L,\qquad I_{N-1} = G.
\]

Por tanto, \(L\) no representa la intensidad del mejor individuo, sino la del
primer individuo no elitista, es decir, del segundo mejor elemento de la
poblaci\'on ordenada.

\subsection{Interpretaci\'on correcta}

Aqu\'i conviene ser muy preciso. RankGA no define por s\'i mismo qu\'e significa
\(I_i\). Solo define \emph{cu\'anto} recibe cada individuo en una escala
abstracta. El significado operacional depende del gen:

\begin{itemize}
  \item el mejor individuo tiene intensidad exactamente \(0\);
  \item el segundo mejor individuo tiene intensidad exactamente \(L\);
  \item el peor individuo tiene intensidad exactamente \(G\).
\end{itemize}

Ese anclaje no es una interpretaci\'on opcional, sino una consecuencia directa
de la forma en que el algoritmo define \(r_i\) y \(\beta\).

\begin{itemize}
  \item en genes reales, \(I_i\) puede ser desviaci\'on est\'andar del ruido;
  \item en genes categ\'oricos, puede usarse como probabilidad de reemplazo;
  \item en otras codificaciones podr\'ia representar radio de vecindad o
  magnitud de perturbaci\'on.
\end{itemize}

Este desacoplamiento hace al framework flexible, pero tambi\'en obliga a que
cada problema mantenga un contrato sem\'antico consistente.

\section{Evaluaci\'on y detecci\'on de mejora}

\subsection{Momento de evaluaci\'on}

Dentro de una iteraci\'on del ciclo, el fitness se recalcula dos veces:

\begin{enumerate}
  \item despu\'es de recombinaci\'on;
  \item despu\'es de mutaci\'on.
\end{enumerate}

La selecci\'on no eval\'ua porque trabaja con copias de individuos cuyo fitness
ya hab\'ia sido conocido antes del clonaje.

\subsection{Criterio de mejora}

El framework actual ya no usa un \'unico criterio de ``mejora'' para todo.
Hay que distinguir entre:

\begin{itemize}
  \item mejora estricta de fitness;
  \item reemplazo del incumbente;
  \item y reinicio de paciencia.
\end{itemize}

En todos los casos se exige adem\'as que el nuevo candidato no sea
genot\'ipicamente id\'entico al incumbente previo. En c\'odigo, la distancia
cuadr\'atica entre ambos debe ser positiva.

\paragraph{Modo \texttt{strict}.}
Si la pol\'itica de incumbente es \texttt{strict}, entonces el reemplazo del
incumbente exige:

\begin{itemize}
  \item el fitness debe ser estrictamente mayor;
  \item la distancia cuadr\'atica entre genomas debe ser positiva.
\end{itemize}

\paragraph{Modo \texttt{neutral}.}
Si la pol\'itica de incumbente es \texttt{neutral}, entonces tambi\'en se
permite sustituir al incumbente cuando:

\begin{itemize}
  \item el fitness es igual al del incumbente actual;
  \item pero la distancia cuadr\'atica entre genomas sigue siendo positiva.
\end{itemize}

Es decir, en modo \texttt{neutral} s\'i se permite que el representante del
mejor fitness alcanzado se desplace sobre una meseta neutral. Lo que no se
permite en ning\'un caso es reemplazarlo por el mismo genotipo.

\subsection{Distancia entre individuos}

La distancia entre individuos se calcula como
\[
d^2(x,y) = \sum_{j=1}^{n} \left( d_j(x_j,y_j) \right)^2,
\]
donde \(d_j\) es la distancia definida por el gen en el locus \(j\).

En genes enteros tipo Hamming unilocus, cada t\'ermino vale \(0\) o \(1\). En
genes reales, depende de la implementaci\'on concreta del gen.

\section{Adaptaci\'on opcional del problema}

Si un problema implementa \texttt{AdaptiveProblem}, entonces el driver llama:
\[
\texttt{adapt(bestFitness)}
\]
despu\'es de completar la iteraci\'on evolutiva.

Esto permite que el problema:

\begin{itemize}
  \item ajuste pesos internos;
  \item cambie par\'ametros de representaci\'on;
  \item modifique penalizaciones;
  \item o altere alguna heur\'istica dependiente del progreso.
\end{itemize}

Si el problema no implementa esa interfaz, no ocurre nada. La adaptaci\'on no
forma parte del contrato b\'asico de \texttt{Problem}.

\section{Criterios de paro}

\subsection{Paro por objetivo}

La corrida termina si el mejor fitness actual satisface
\[
f_{\max} \geq f_{\text{goal}},
\]
donde \(f_{\text{goal}} = \texttt{problem.getGoalFt()}\).

\subsection{Paro por paciencia}

Tambi\'en termina si transcurre un tiempo fijo sin el tipo de progreso que la
pol\'itica de paciencia considere suficiente.

La versi\'on actual permite parametrizar dos cosas distintas:

\begin{itemize}
  \item el criterio de reemplazo del incumbente;
  \item y el criterio de reinicio de paciencia.
\end{itemize}

La paciencia se configura con un par\'ametro en milisegundos:
\[
\texttt{patienceMillis}.
\]

Por defecto:
\[
\texttt{patienceMillis} = 60\,000 \text{ ms} = 1 \text{ minuto}.
\]

\subsection{Pol\'itica de incumbente}

El framework distingue dos modos:

\begin{itemize}
  \item \texttt{strict}: el incumbente solo se sustituye si aparece un
  individuo con fitness estrictamente mayor;
  \item \texttt{neutral}: el incumbente tambi\'en puede sustituirse por otro
  genotipo con el mismo fitness, siempre que la distancia genot\'ipica sea
  positiva.
\end{itemize}

Esto controla si el representante del mejor fitness alcanzado puede o no
desplazarse sobre una meseta neutral.

\subsection{Pol\'itica de reinicio de paciencia}

El reinicio de paciencia tambi\'en es configurable:

\begin{itemize}
  \item \texttt{fitness}: la paciencia se reinicia solo con mejora estricta de
  fitness;
  \item \texttt{movement}: la paciencia se reinicia cuando cambia el
  incumbente seg\'un la pol\'itica de incumbente seleccionada.
\end{itemize}

La distinci\'on es importante. Permitir movimiento neutral del incumbente no
obliga a interpretar ese movimiento como progreso suficiente para seguir
extendiendo la corrida. Son dos decisiones metodol\'ogicas distintas.

\subsection{Consecuencia importante}

La paciencia es \textbf{por repetici\'on}, no por experimento completo. Si se
ejecutan \(100\) repeticiones y ninguna alcanza el objetivo, el tiempo total
puede escalar aproximadamente a \(100\) minutos, salvo peque\~nas variaciones
por inicializaci\'on, I/O y sistema operativo.

\section{Contabilidad de evaluaciones}

\subsection{Qu\'e se cuenta}

La variable \texttt{evaluations} cuenta cu\'antas veces se evalu\'o fitness en
una repetici\'on.

Se suma:

\begin{itemize}
  \item \(N\) al inicializar;
  \item \(N\) despu\'es de recombinaci\'on en cada generaci\'on;
  \item \(N\) despu\'es de mutaci\'on en cada generaci\'on.
\end{itemize}

\subsection{Correcci\'on importante ya incorporada}

La versi\'on actual ya no fuerza una generaci\'on extra si la poblaci\'on
inicial alcanza el objetivo. El ciclo principal usa un \texttt{while} y no un
\texttt{do-while}. Eso evita sobrecontar evaluaciones en casos donde la
soluci\'on objetivo ya aparece desde el principio.

\section{Salida experimental}

\subsection{Logs heredados}

Por cada corrida se generan:

\begin{itemize}
  \item un archivo \texttt{.txt} con snapshots cronol\'ogicos del mejor
  individuo;
  \item un archivo \texttt{\_rep.txt} por repetici\'on con el \'ultimo estado de
  la poblaci\'on en esa repetici\'on.
\end{itemize}

\subsection{Resumen estructurado}

Tambi\'en se genera un \texttt{summary.csv} con una fila por repetici\'on. Entre
sus columnas est\'an:

\begin{itemize}
  \item algoritmo;
  \item clase del problema;
  \item nombre del problema;
  \item par\'ametros efectivos del problema;
  \item semilla;
  \item \texttt{run\_id};
  \item \'indice de repetici\'on;
  \item tama\~no de poblaci\'on;
  \item n\'umero total de repeticiones;
  \item n\'umero de evaluaciones;
  \item mejor fitness final;
  \item tiempo transcurrido en milisegundos;
  \item raz\'on de terminaci\'on;
  \item fitness objetivo;
  \item paciencia configurada en milisegundos;
  \item pol\'itica de reemplazo del incumbente;
  \item pol\'itica de reinicio de paciencia;
  \item prefijo del archivo de salida.
\end{itemize}

\subsection{Gr\'aficas autom\'aticas}

Al terminar todas las repeticiones, RankGA invoca un script en Python que
genera:

\begin{itemize}
  \item una gr\'afica ordenada por evaluaciones;
  \item una gr\'afica ordenada por tiempo transcurrido;
  \item y dos CSV ordenados correspondientes.
\end{itemize}

La versi\'on actual ya incorpora el \texttt{runId} en el nombre de la figura
para evitar colisiones entre corridas con el mismo problema y la misma semilla.

\section{Ejemplo: OneMax de 8 bits}

\subsection{Definici\'on}

En el problema OneMax de longitud \(n=8\):

\begin{itemize}
  \item el genoma es binario;
  \item el fitness es la suma de bits en uno;
  \item el objetivo es alcanzar fitness \(8\).
\end{itemize}

Formalmente, si \(x \in \{0,1\}^8\), entonces
\[
f(x) = \sum_{j=1}^{8} x_j.
\]

\subsection{Configuraci\'on en el repositorio}

La implementaci\'on actual:

\begin{itemize}
  \item usa \texttt{GeneInteger} con dominio binario;
  \item define intensidad local \(L = 1/8\);
  \item define intensidad global \(G = 0.5\);
  \item usa una estrategia de copia de genes que preserva la independencia
  estoc\'astica entre individuos.
\end{itemize}

\subsection{Interpretaci\'on experimental}

OneMax es un benchmark deliberadamente sencillo. Si este problema falla con
frecuencia por paciencia, eso suele indicar:

\begin{itemize}
  \item un defecto en la implementaci\'on;
  \item una sem\'antica inconsistente de mutaci\'on;
  \item o un error en el dise\~no experimental y no una limitaci\'on genuina
  del algoritmo.
\end{itemize}

\section{Supuestos y limitaciones del framework}

\subsection{Supuestos fuertes}

El algoritmo presupone:

\begin{itemize}
  \item fitness orientado a maximizaci\'on;
  \item poblaciones de tama\~no al menos \(3\);
  \item \(G > L\) para que la mutaci\'on dependa del rango de forma mon\'otona;
  \item una implementaci\'on consistente del operador de copia de genes;
  \item y una interpretaci\'on bien definida de intensidad dentro de cada tipo
  de gen.
\end{itemize}

\subsection{Decisiones de dise\~no}

\begin{itemize}
  \item La presi\'on selectiva \(S\) es fija.
  \item El emparejamiento para recombinaci\'on es adyacente y no aleatorio.
  \item La selecci\'on usa redondeo estoc\'astico secuencial y no necesita
  muestreo multinomial exacto para cumplir su objetivo.
\end{itemize}

\subsection{Limitaciones reales de implementaci\'on}

\begin{itemize}
  \item El motor usa estado est\'atico en \texttt{RankGA}, por lo que est\'a
  dise\~nado como \emph{runner} simple, no como motor concurrente para varias
  corridas simult\'aneas en la misma JVM.
  \item La generaci\'on autom\'atica de gr\'aficas depende de componentes
  externos disponibles en el entorno.
\end{itemize}

\subsection{Criterio documental}

Para entender el comportamiento real del framework, hay que preferir:

\begin{enumerate}
  \item el c\'odigo ejecutable;
  \item la ejecuci\'on efectiva de los tests autom\'aticos del proyecto, cuando
  son aplicables al cambio o a la afirmaci\'on que se quiere sostener;
  \item las corridas reproducibles y sus res\'umenes experimentales
  estructurados;
  \item y la revisi\'on manual de archivos de salida cuando hace falta una
  validaci\'on complementaria.
\end{enumerate}

\subsection{Observaci\'on pr\'actica sobre verificaci\'on}

En este repositorio conviene distinguir entre:

\begin{itemize}
  \item compilaci\'on correcta del c\'odigo;
  \item ejecuci\'on efectiva de pruebas autom\'aticas;
  \item y validaci\'on manual mediante corridas, inspecci\'on de logs,
  \texttt{summary.csv} y gr\'aficas.
\end{itemize}

Desde el punto de vista metodol\'ogico, esa distinci\'on importa porque una
afirmaci\'on del tipo ``esto ya est\'a probado'' no significa lo mismo si se
basa en una bater\'ia automatizada completa, en una corrida reproducible o en
una inspecci\'on manual del \emph{summary}. En una descripci\'on rigurosa del
framework, esas fuentes de evidencia no deben confundirse ni presentarse como
si tuvieran exactamente el mismo alcance.

\section{Resumen conceptual final}

RankGA puede resumirse como sigue:

\begin{enumerate}
  \item ordena la poblaci\'on por fitness;
  \item replica individuos seg\'un su rango, no seg\'un el valor crudo del
  fitness;
  \item cruza por parejas adyacentes;
  \item muta con intensidad creciente conforme empeora el rango;
  \item preserva al mejor individuo porque su intensidad es cero;
  \item asigna al segundo mejor individuo exactamente la intensidad local
  \(L\);
  \item permite adaptaci\'on opcional del problema;
  \item y produce salidas estructuradas para an\'alisis experimental.
\end{enumerate}

La contribuci\'on pr\'actica del framework no es solo combinar selecci\'on por
rango con mutaci\'on dependiente del rango, sino hacerlo de manera gen\'erica
para varias codificaciones y problemas. El precio de esa generalidad es que el
contrato entre motor y genes debe estar especialmente bien cuidado. Cuando ese
contrato se respeta, el comportamiento es razonable y reproducible; cuando se
viola, incluso problemas muy f\'aciles pueden dar resultados enga\~nosos.

\end{document}
